\documentclass[letterpaper,11pt]{article}

\usepackage[utf8]{inputenc}
\usepackage[T1]{fontenc}
\usepackage[turkish]{babel}

\usepackage{latexsym}
\usepackage[empty]{fullpage}
\usepackage{titlesec}
\usepackage{marvosym}
\usepackage[usenames,dvipsnames]{color}
\usepackage{verbatim}
\usepackage{enumitem}
\usepackage[hidelinks]{hyperref}
\usepackage{fancyhdr}
\usepackage{tabularx}
\usepackage[sfdefault]{FiraSans}
\input{glyphtounicode}

\pagestyle{fancy}
\fancyhf{} 
\fancyfoot{}
\renewcommand{\headrulewidth}{0pt}
\renewcommand{\footrulewidth}{0pt}

\addtolength{\oddsidemargin}{-0.5in}
\addtolength{\evensidemargin}{-0.5in}
\addtolength{\textwidth}{1in}
\addtolength{\topmargin}{-.5in}
\addtolength{\textheight}{1.0in}

\urlstyle{same}

\raggedbottom
\raggedright
\setlength{\tabcolsep}{0in}

\titleformat{\section}{
  \vspace{10pt}\scshape\raggedright\large
}{}{0em}{}[\color{black}\titlerule \vspace{-5pt}]

\pdfgentounicode=1

\newcommand{\resumeItem}[1]{
  \item\small{
    {#1 \vspace{-2pt}}
  }
}

\newcommand{\resumeSubheading}[4]{
  \vspace{-2pt}\item
    \begin{tabular*}{0.97\textwidth}[t]{l@{\extracolsep{\fill}}r}
      \textbf{#1} & \small #2 \\
      \textit{\small#3} & \textit{\small #4} \\
    \end{tabular*}\vspace{-7pt}
}

\newcommand{\resumeProjectHeading}[2]{
    \item
    \begin{tabular*}{0.97\textwidth}{l@{\extracolsep{\fill}}r}
      \small#1 & \small #2 \\
    \end{tabular*}\vspace{-7pt}
}

\newcommand{\resumeSubHeadingListStart}{\begin{itemize}[leftmargin=0.15in, label={}]}
\newcommand{\resumeSubHeadingListEnd}{\end{itemize}}
\newcommand{\resumeItemListStart}{\begin{itemize}}
\newcommand{\resumeItemListEnd}{\end{itemize}\vspace{-5pt}}

\begin{document}

\begin{center}
    \textbf{\Huge \scshape Ömer Gülçiçek} \\ \vspace{1pt}
    \small 0554 350 98 06 $|$ \href{mailto:iletisim@omergulcicek.com}{\underline{iletisim@omergulcicek.com}} $|$ \href{https://omergulcicek.com/}{\underline{omergulcicek.com}} \\
    \small \href{https://linkedin.com/in/omergulcicek}{\underline{linkedin.com/in/omergulcicek}} $|$ \href{https://github.com/omergulcicek}{\underline{github.com/omergulcicek}} \\
    \small İstanbul, Gaziosmanpaşa $|$ Askerlik: Yapıldı
\end{center}

\section{Özet Bilgi}
\small{
React ve Next.js ekosisteminde 8 yılı aşkın tecrübeyle; hem kullanıcılar hem de geliştiriciler için sürdürülebilir ürünler inşa ediyorum. AI-native kural setleri geliştirerek, ekiplerin DX standartlarını ve mühendislik verimliliğini artırmaya odaklanıyorum. Büyük ölçekli frontend projelerinde performansı artırmaya, takıma yol göstermeye ve işi verimli kılan yazılım araçları geliştirmeye odaklıyım.
}

\section{Teknik Uzmanlık}
 \begin{itemize}[leftmargin=0.15in, label={}]
    \small{\item{
     \textbf{Frontend Development}{: React, Next.js, TypeScript, JavaScript (ES6+), TanStack (Start, Query, Router)} \\
     \textbf{State \& Form Management}{: Zustand, Zod, React Hook Form} \\
     \textbf{Styling \& UI}{: Tailwind CSS, shadcn/ui, Framer Motion, Storybook, CSS3/SASS, HTML5} \\
     \textbf{Araçlar \& Veri}{: Git, Cursor (AI-Native Development), Figma, GraphQL} \\
     \textbf{Diller}{: Türkçe (Anadil), İngilizce (Teknik okuma ve yazma)}
    }}
 \end{itemize}

\section{İş Tecrübeleri}
  \resumeSubHeadingListStart

    \resumeSubheading
      {Frontend Team Lead}{Ağu 2025 -- Kas 2025}
      {Machinarium}{Remote}
      \resumeItemListStart
        \resumeItem{Yapay zekâ kuralları ve kod yazım standartları oluşturarak geliştirici verimliliğini ve kod tutarlılığını artırdım.}
        \resumeItem{10 kişilik frontend ekibine teknik rehberlik ederek kompleks ürünlerin yüksek mühendislik standartlarında yayına alınmasını sağladım.}
        \resumeItem{Ortak bileşen kullanımı ve modüler yapıları kurgulayarak e-ticaret ekosistemindeki (PIM, CMS, OMS) teknik tutarsızlıkları giderdim.}
      \resumeItemListEnd

    \resumeSubheading
      {Senior Frontend Developer}{Mar 2024 -- Haz 2025}
      {Alışgidiş}{Remote}
      \resumeItemListStart
        \resumeItem{alisgidis.com'un ölçeklenebilir mimari temellerini inşa ettim; düzenli teknik oturumlar ve mentorlukla ekibin mühendislik standartlarını yükselttim.}
        \resumeItem{Core Web Vitals iyileştirmeleri ve kod optimizasyonlarıyla Lighthouse skorlarını 90+ seviyesine taşıyarak site hızını ve SEO verimliliğini artırdım.}
        \resumeItem{Atomic Design prensipleriyle modüler ve dark/light tema destekli bir UI kütüphanesi kurgulayarak frontend geliştirme hızını \%25 artırdım.}
      \resumeItemListEnd

    \resumeSubheading
      {Senior Frontend Developer}{Oca 2022 -- Şub 2024}
      {Gordion}{Remote}
      \resumeItemListStart
        \resumeItem{Modern mühendislik pratiklerini ve kod kalitesi standartlarını ekip kültürüne entegre ederek teknik borcu minimize ettim.}
        \resumeItem{Kurumsal UI Kit ve tasarım sistemleri geliştirerek ekipler arası teknik tutarlılığı sağladım ve geliştirme hızını artırdım.}
        \resumeItem{Kapsamlı rezervasyon platformları için ölçeklenebilir frontend mimarileri inşa ettim; GraphQL entegrasyonu ile veri yönetim süreçlerini optimize ettim.}
      \resumeItemListEnd

    \resumeSubheading
      {Frontend Developer}{Eyl 2020 -- Oca 2022}
      {HangiKredi}{Ofisten}
      \resumeItemListStart
        \resumeItem{hangikredi.com platformunu uçtan uca geliştirerek performans iyileştirmelerini yönettim; yüksek performanslı arayüzler ve yeniden kullanılabilir kütüphanelerle ekip standartlarını güçlendirdim.}
        \resumeItem{Figma üzerindeki kompleks tasarımları yüksek hassasiyetle koda dönüştürerek ürün geliştirme sürecini hızlandırdım.}
      \resumeItemListEnd
      
    \vspace{30pt}

    \resumeSubheading
      {Frontend Eğitmeni}{Ara 2020 -- Şub 2021}
      {Kodluyoruz}{Remote}
      \resumeItemListStart
        \resumeItem{30 öğrenciye HTML, CSS, SASS, JavaScript, Git ve React konularını kapsayan kapsamlı Frontend eğitimi verdim.}
        \resumeItem{PSD to CSS pratikleri ve teknik mentorluk süreçleriyle öğrencilerin modern arayüz geliştirme becerilerini artırdım.}
      \resumeItemListEnd

    \resumeSubheading
      {Frontend Developer}{Haz 2019 -- Kas 2019}
      {Akinon}{Ofisten}
      \resumeItemListStart
        \resumeItem{Vakko ve A101 gibi ölçekli e-ticaret markaları için modüler bileşen mimarileri kurguladım; Zeplin tasarımlarını yüksek sadakatle hayata geçirdim.}
        \resumeItem{Hazırladığım teknik dokümantasyonlar ile geliştirme süreçlerinin sürdürülebilirliğini ve ekip içi bilgi aktarımını sağladım.}
      \resumeItemListEnd

    \resumeSubheading
      {Frontend Developer}{Tem 2018 -- Kas 2018}
      {Turkcell}{Ofisten}
      \resumeItemListStart
        \resumeItem{Kritik frontend problemlerine gerçek zamanlı çözümler üreterek geliştirme döngüsündeki blokajları giderdim.}
        \resumeItem{Canlı sistemdeki hataları anlık olarak optimize ederek ekipler arası geri bildirim hızını ve sistem kararlılığını artırdım.}
      \resumeItemListEnd

    \resumeSubheading
      {Frontend Developer}{Tem 2017 -- Haz 2018}
      {Ziraat Teknoloji}{Ofisten}
      \resumeItemListStart
        \resumeItem{İnşaat projelerinin bütçe ve tarih takibinin yapıldığı kapsamlı yönetim panelini teknik standartlara uygun şekilde hayata geçirdim.}
        \resumeItem{Kendo UI tabanlı kurumsal admin paneli projelerini uçtan uca geliştirip yönettim; arayüzleri responsive hale getirerek kullanıcı deneyimini modernize ettim.}
      \resumeItemListEnd

  \resumeSubHeadingListEnd
  
\section{Eğitim Bilgileri}
  \resumeSubHeadingListStart
    \resumeSubheading
      {Yazılım Mühendisliği, İngilizce}{2012 -- 2017}
      {İstanbul Aydın Üniversitesi}{İstanbul}
  \resumeSubHeadingListEnd

\section{Projeler ve Açık Kaynak}
    \resumeSubHeadingListStart
      \resumeProjectHeading
          {\textbf{ViraStack AI} $|$ \href{https://github.com/virastack/ai-rules}{\underline{github.com/virastack/ai-rules}}}{2026}
          \resumeItemListStart
            \resumeItem{Modern React uygulamaları için yapay zeka odaklı mimari kiti ve yüksek disiplinli protokoller.}
          \resumeItemListEnd

      \resumeProjectHeading
          {\textbf{Input Mask} $|$ \href{https://github.com/virastack/input-mask}{\underline{github.com/virastack/input-mask}}}{2025}
          \resumeItemListStart
            \resumeItem{React Hook Form için optimize edilmiş, hafif ve bağımlılıksız input maskeleme kütüphanesi.}
          \resumeItemListEnd

      \resumeProjectHeading
          {\textbf{Password Toggle} $|$ \href{https://github.com/virastack/password-toggle}{\underline{github.com/virastack/password-toggle}}}{2025}
          \resumeItemListStart
            \resumeItem{React için tam erişilebilir ve son derece özelleştirilebilir şifre görünürlük hook'u.}
          \resumeItemListEnd
          
      \resumeProjectHeading
          {\textbf{Next.js Boilerplate} $|$ \href{https://github.com/virastack/nextjs-boilerplate}{\underline{github.com/virastack/nextjs-boilerplate}}}{2024}
          \resumeItemListStart
            \resumeItem{Tailwind CSS 4 ve TypeScript ile oluşturulmuş, üretime hazır Next.js 16+ başlangıç şablonu.}
          \resumeItemListEnd
          
      \resumeProjectHeading
          {\textbf{Rûznâme} $|$ \href{https://github.com/omergulcicek/ruzname}{\underline{github.com/omergulcicek/ruzname}}}{2025}
          \resumeItemListStart
            \resumeItem{Takvim abonelikleri için ücretsiz, açık kaynak takvim arayüzleri ve içerik platformu.}
          \resumeItemListEnd

      \resumeProjectHeading
          {\textbf{Türkçe Doküman} $|$ \href{https://github.com/omergulcicek/turkcedokuman.com}{\underline{github.com/omergulcicek/turkcedokuman.com}}}{2023}
          \resumeItemListStart
            \resumeItem{Yazılım geliştirme konularında Türkçe kaynak ve içerik projesi.}
          \resumeItemListEnd

      \resumeProjectHeading
        {\textbf{Turkuaz} $|$ \href{https://github.com/omergulcicek/turkuaz}{\underline{github.com/omergulcicek/turkuaz}}}{2021}
        \resumeItemListStart
          \resumeItem{Hızlı ve minimalist web sayfaları için özelleşmiş CSS kütüphanesi.}
        \resumeItemListEnd
    \resumeSubHeadingListEnd
    
    \section{Kişisel İlgi Alanları}
         \small{
          Teknik dünyamın dışında; kişisel blogumda teknik konuların yanı sıra deneyimlerimi paylaştığım \textbf{yazılar} yazıyorum. Düzenli \textbf{satranç} oynuyor, doğa yürüyüşleri yapıp \textbf{fotoğraf} çekiyorum. \textbf{Futbol}, \textbf{masa tenisi}, \textbf{bisiklet} ve \textbf{yüzme} sporlarıyla ilgileniyorum. Tarih (Osmanlı-Roma), \textbf{dini kaynaklar} ve kişisel gelişim üzerine \textbf{kitap} okumayı seviyorum.
         }
\end{document}
